\documentclass[11pt,a4paper]{article}

% Packages
\usepackage[utf8]{inputenc}
\usepackage[T1]{fontenc}
\usepackage{lmodern}
\usepackage[margin=1in]{geometry}
\usepackage{graphicx}
\usepackage{hyperref}
\usepackage{xcolor}
\usepackage{listings}
\usepackage{booktabs}
\usepackage{amsmath}
\usepackage{amssymb}
\usepackage{fancyhdr}
\usepackage{titlesec}
\usepackage{tocloft}
\usepackage{enumitem}
\usepackage{longtable}

% Colors
\definecolor{codegreen}{rgb}{0,0.6,0}
\definecolor{codegray}{rgb}{0.5,0.5,0.5}
\definecolor{codepurple}{rgb}{0.58,0,0.82}
\definecolor{backcolour}{rgb}{0.95,0.95,0.92}
\definecolor{linkblue}{rgb}{0.0,0.4,0.7}

% Hyperref setup
\hypersetup{
    colorlinks=true,
    linkcolor=linkblue,
    filecolor=magenta,
    urlcolor=linkblue,
    pdftitle={Intensity Measure - ImageJ ROI Intensity Analysis},
    pdfauthor={},
}

% Code listing style
\lstdefinestyle{mystyle}{
    backgroundcolor=\color{backcolour},
    commentstyle=\color{codegreen},
    keywordstyle=\color{magenta},
    numberstyle=\tiny\color{codegray},
    stringstyle=\color{codepurple},
    basicstyle=\ttfamily\footnotesize,
    breakatwhitespace=false,
    breaklines=true,
    captionpos=b,
    keepspaces=true,
    numbers=left,
    numbersep=5pt,
    showspaces=false,
    showstringspaces=false,
    showtabs=false,
    tabsize=2,
    frame=single,
    framerule=0pt,
}
\lstset{style=mystyle}

% Header/Footer
\pagestyle{fancy}
\fancyhf{}
\rhead{Intensity Measure Documentation}
\lhead{ImageJ ROI Intensity Analysis}
\rfoot{Page \thepage}

% Title formatting
\titleformat{\section}
    {\normalfont\Large\bfseries\color{linkblue}}{\thesection}{1em}{}
\titleformat{\subsection}
    {\normalfont\large\bfseries}{\thesubsection}{1em}{}
\titleformat{\subsubsection}
    {\normalfont\normalsize\bfseries}{\thesubsubsection}{1em}{}

% Document info
\title{
    \vspace{-1cm}
    \textbf{\LARGE Intensity Measure}\\[0.5cm]
    \large ImageJ ROI Intensity Analysis Script\\[0.3cm]
    \normalsize Documentation \& User Guide\\[0.2cm]
    \small Version 2.1
}
\author{Radmir Sarsenov}
\date{\today}

\begin{document}

\maketitle
\thispagestyle{empty}

\vspace{1cm}

\begin{abstract}
\noindent
A JavaScript macro for ImageJ/Fiji that measures fluorescence intensity across image stacks with \textbf{multiple normalization methods}. Suitable for time-lapse imaging, photobleaching studies, calcium imaging, FRAP, and general fluorescence quantification. This script offers six different normalization methods including F/F$_0$, $\Delta$F/F$_0$, Z-score, and more.
\end{abstract}

\vspace{0.5cm}
\tableofcontents
\newpage

%==============================================================================
\section{Overview}
%==============================================================================

This script automates the process of measuring fluorescence intensity changes across a time-lapse image stack. It offers \textbf{six different normalization methods} to suit various experimental needs, from photobleaching analysis to calcium imaging to protein expression quantification.

\subsection{Use Cases}

\begin{itemize}[leftmargin=*]
    \item \textbf{Photobleaching / FRAP}: Track fluorescence decay or recovery over time
    \item \textbf{Calcium imaging}: Monitor intracellular calcium dynamics
    \item \textbf{Drug response}: Measure intensity changes upon treatment
    \item \textbf{Expression analysis}: Compare fluorescence levels across conditions
    \item \textbf{Time-lapse studies}: Quantify any temporal intensity changes
\end{itemize}

%==============================================================================
\section{Features}
%==============================================================================

\begin{itemize}[leftmargin=*]
    \item \textbf{Six normalization methods}: F/F$_0$, $\Delta$F/F$_0$, background subtraction, percent of max, Z-score, or raw values
    \item \textbf{Flexible background selection}: Click on the image, type coordinates, or skip background correction
    \item \textbf{Multiple ROI tools}: Rectangle, Oval, Freehand, Polygon, Brush, or use existing ROI
    \item \textbf{Per-frame background correction}: Accounts for fluctuating background
    \item \textbf{CTCF calculation}: Corrected Total Cell Fluorescence for every frame
    \item \textbf{Validation checks}: Ensures selections are within image bounds
    \item \textbf{Results table output}: Easy export to CSV for further analysis
\end{itemize}

%==============================================================================
\section{Installation}
%==============================================================================

\subsection{Requirements}

\begin{itemize}
    \item \href{https://imagej.nih.gov/ij/}{ImageJ} (v1.53+) or \href{https://fiji.sc/}{Fiji}
\end{itemize}

\subsection{Setup}

\begin{enumerate}
    \item Download \texttt{Intensity\_measure.js}
    \item Place it in your ImageJ plugins folder:
    \begin{itemize}
        \item \textbf{Windows}: \texttt{C:\textbackslash Program Files\textbackslash ImageJ\textbackslash plugins\textbackslash}
        \item \textbf{macOS}: \texttt{/Applications/ImageJ.app/plugins/}
        \item \textbf{Linux}: \texttt{\textasciitilde/ImageJ/plugins/}
    \end{itemize}
    \item Restart ImageJ, or use \texttt{Plugins > Macros > Run...} to execute directly
\end{enumerate}

%==============================================================================
\section{Usage}
%==============================================================================

\subsection{Step-by-Step Guide}

\subsubsection{Step 1: Open Your Image Stack}

Open your time-lapse fluorescence image stack in ImageJ:
\begin{itemize}
    \item \texttt{File > Open...} or drag-and-drop
    \item Ensure it's a stack (multiple frames), not a single image
\end{itemize}

\subsubsection{Step 2: Run the Script}

\begin{itemize}
    \item \texttt{Plugins > Macros > Run...} $\rightarrow$ Select \texttt{Intensity\_measure.js}
    \item Or if installed: \texttt{Plugins > Intensity\_measure}
\end{itemize}

\subsubsection{Step 3: Select Normalization Method}

Choose from six normalization options:

\begin{table}[h]
\centering
\begin{tabular}{@{}lll@{}}
\toprule
\textbf{Method} & \textbf{Output Range} & \textbf{Best For} \\
\midrule
F/F$_0$ (default) & Starts at 1.0 & Photobleaching, FRAP, general time-lapse \\
$\Delta$F/F$_0$ & Starts at 0.0 & Calcium imaging, stimulus response \\
Background subtraction & Absolute units & Simple correction, comparing intensities \\
Percent of maximum & 0--100\% & Normalizing across experiments \\
Z-score & Mean=0, SD=1 & Statistical comparison, outlier detection \\
Raw values & Original units & No processing needed \\
\bottomrule
\end{tabular}
\caption{Available normalization methods}
\end{table}

\subsubsection{Step 4: Select Background Point}

You'll see a dialog asking how to select the background reference:

\begin{table}[h]
\centering
\begin{tabular}{@{}ll@{}}
\toprule
\textbf{Option} & \textbf{Description} \\
\midrule
Click to select & Use the Point tool to click on a background region \\
Type coordinates & Enter exact X, Y pixel coordinates \\
No background correction & Skip background subtraction entirely \\
\bottomrule
\end{tabular}
\end{table}

\textbf{Tips for background selection:}
\begin{itemize}
    \item Choose a region with NO fluorescent signal
    \item Avoid dark artifacts or dead pixels
    \item The same pixel location is used across all frames
\end{itemize}

\subsubsection{Step 5: Select Target ROI}

Choose your preferred ROI selection method:

\begin{table}[h]
\centering
\begin{tabular}{@{}ll@{}}
\toprule
\textbf{Option} & \textbf{Description} \\
\midrule
Rectangle & Draw a rectangular selection \\
Oval & Draw an elliptical selection \\
Freehand & Draw a freeform closed shape \\
Polygon & Click vertices to create a polygon \\
Brush & Paint the selection with a brush \\
Use existing ROI & Use a previously drawn ROI on the image \\
Type coordinates (Rectangle) & Enter X, Y, Width, Height for a rectangle \\
Type coordinates (Oval) & Enter X, Y, Width, Height for an oval \\
\bottomrule
\end{tabular}
\end{table}

\subsubsection{Step 6: View Results}

A Results Table will appear with measurements for each slice:

\begin{lstlisting}[language={},numbers=none]
Slice | Raw_Mean | Background | F_corrected | F/F0   | CTCF
------+----------+------------+-------------+--------+--------
1     | 1500.0   | 100.0      | 1400.0      | 1.0000 | 56000.0
2     | 1450.0   | 102.0      | 1348.0      | 0.9629 | 53920.0
3     | 1380.0   | 99.0       | 1281.0      | 0.9150 | 51240.0
\end{lstlisting}

\subsubsection{Step 7: Export Data}

\begin{itemize}
    \item \texttt{File > Save As...} in the Results window
    \item Save as \texttt{.csv} for Excel, Python, R, MATLAB, etc.
\end{itemize}

%==============================================================================
\section{Normalization Methods}
%==============================================================================

\subsection{F/F$_0$ (Normalize to First Frame)}

\textbf{Formula:}
\begin{equation}
\frac{F}{F_0} = \frac{F(t)}{F(t_0)}
\end{equation}

\begin{itemize}
    \item Output starts at \textbf{1.0} and changes from there
    \item Standard for \textbf{photobleaching}, \textbf{FRAP}, general time-lapse
    \item Easy to interpret: 0.5 = 50\% of initial intensity
\end{itemize}

\subsection{$\Delta$F/F$_0$ (Fold Change from Baseline)}

\textbf{Formula:}
\begin{equation}
\frac{\Delta F}{F_0} = \frac{F(t) - F(t_0)}{F(t_0)}
\end{equation}

\begin{itemize}
    \item Output starts at \textbf{0.0} (baseline)
    \item Standard for \textbf{calcium imaging}, stimulus-response experiments
    \item Positive values = increase, negative = decrease
    \item A value of 1.0 means 100\% increase (doubled)
\end{itemize}

\subsection{Background Subtraction Only}

\textbf{Formula:}
\begin{equation}
F_{\text{corrected}} = F(t) - F_{\text{background}}(t)
\end{equation}

\begin{itemize}
    \item Returns absolute intensity values (background-corrected)
    \item Good for comparing intensities within the same experiment
    \item No normalization to baseline
\end{itemize}

\subsection{Percent of Maximum}

\textbf{Formula:}
\begin{equation}
\% = \frac{F - F_{\min}}{F_{\max} - F_{\min}} \times 100
\end{equation}

\begin{itemize}
    \item Scales data to \textbf{0--100\%} range
    \item Useful for comparing across different experiments
    \item Min intensity = 0\%, Max intensity = 100\%
\end{itemize}

\subsection{Z-score Normalization}

\textbf{Formula:}
\begin{equation}
Z = \frac{F - \mu}{\sigma}
\end{equation}

\begin{itemize}
    \item Mean becomes 0, standard deviation becomes 1
    \item Good for \textbf{statistical comparisons} across experiments
    \item Identifies outliers ($|Z| > 2$ or 3)
\end{itemize}

\subsection{Raw Values (No Normalization)}

\begin{itemize}
    \item Returns the original mean intensity values
    \item No background correction or normalization applied
    \item Use when you need unprocessed data
\end{itemize}

%==============================================================================
\section{Output Columns Explained}
%==============================================================================

\begin{table}[h]
\centering
\begin{tabular}{@{}p{2.5cm}p{4.5cm}p{6cm}@{}}
\toprule
\textbf{Column} & \textbf{Formula} & \textbf{Description} \\
\midrule
Slice & --- & Frame number (1-indexed) \\
Raw\_Mean & $\text{mean}(\text{ROI pixels})$ & Uncorrected mean intensity within ROI \\
Background & $\text{pixel}(x_{ref}, y_{ref})$ & Background value at reference point (if used) \\
F\_corrected & $\text{Raw\_Mean} - \text{Background}$ & Background-subtracted mean intensity \\
[Normalized] & Varies by method & The normalized value (column name depends on method) \\
CTCF & $(A \times \bar{I}) - (A \times \bar{I}_{bg})$ & Corrected Total Cell Fluorescence \\
\bottomrule
\end{tabular}
\end{table}

%==============================================================================
\section{Mathematical Background}
%==============================================================================

\subsection{F/F$_0$ Normalization}

The \textbf{F/F$_0$} ratio normalizes all measurements to the first frame:

\begin{equation}
\frac{F}{F_0} = \frac{F(t) - F_{bg}(t)}{F(t_0) - F_{bg}(t_0)}
\end{equation}

Where:
\begin{itemize}
    \item $F(t)$ = Mean fluorescence intensity at time $t$
    \item $F_{bg}(t)$ = Background intensity at time $t$
    \item $F(t_0)$ = Mean fluorescence at the first frame
    \item $F_{bg}(t_0)$ = Background at the first frame
\end{itemize}

\subsection{$\Delta$F/F$_0$ (Fold Change)}

\begin{equation}
\frac{\Delta F}{F_0} = \frac{F(t) - F_0}{F_0} = \frac{F}{F_0} - 1
\end{equation}

This is simply F/F$_0$ minus 1, shifting the baseline to 0.

\subsection{CTCF (Corrected Total Cell Fluorescence)}

CTCF accounts for both intensity and area:

\begin{equation}
\text{CTCF} = (A_{\text{ROI}} \times \bar{I}_{\text{ROI}}) - (A_{\text{ROI}} \times \bar{I}_{\text{background}})
\end{equation}

This is useful when comparing absolute fluorescence amounts across different sized structures.

\subsection{Exponential Decay Fitting}

For photobleaching analysis, fit your F/F$_0$ data to:

\begin{equation}
\frac{F}{F_0} = A \cdot e^{-t/\tau} + B
\end{equation}

Where:
\begin{itemize}
    \item $\tau$ = photobleaching time constant
    \item $A$ = bleachable fraction
    \item $B$ = non-bleachable baseline
\end{itemize}

\textbf{Half-life calculation:}
\begin{equation}
t_{1/2} = \tau \cdot \ln(2) \approx 0.693 \cdot \tau
\end{equation}

%==============================================================================
\section{Best Practices}
%==============================================================================

\subsection{For Time-Lapse Experiments}

\begin{enumerate}
    \item \textbf{Consistent imaging conditions}
    \begin{itemize}
        \item Same exposure time for all frames
        \item Same excitation intensity
        \item Let laser/lamp stabilize before imaging
    \end{itemize}
    
    \item \textbf{Background selection}
    \begin{itemize}
        \item Choose a region that never contains signal
        \item Avoid edges of the field of view
        \item Check that background doesn't change dramatically
    \end{itemize}
    
    \item \textbf{ROI selection}
    \begin{itemize}
        \item Include the entire fluorescent structure
        \item Don't include neighboring cells/structures
        \item Use the same ROI for all comparisons
    \end{itemize}
\end{enumerate}

\subsection{Choosing a Normalization Method}

\begin{table}[h]
\centering
\begin{tabular}{@{}ll@{}}
\toprule
\textbf{Experiment Type} & \textbf{Recommended Method} \\
\midrule
Photobleaching & F/F$_0$ \\
FRAP & F/F$_0$ \\
Calcium imaging & $\Delta$F/F$_0$ \\
Drug response & $\Delta$F/F$_0$ or F/F$_0$ \\
Comparing expression levels & Background subtraction or CTCF \\
Cross-experiment comparison & Percent of max or Z-score \\
Raw data export & Raw values \\
\bottomrule
\end{tabular}
\caption{Recommended normalization methods by experiment type}
\end{table}

%==============================================================================
\section{Troubleshooting}
%==============================================================================

\subsection{``Initial ROI intensity is <= background''}

\textbf{Cause:} Your selected ROI has lower intensity than the background reference point.

\textbf{Solution:}
\begin{itemize}
    \item Check that the background point is truly in a dark region
    \item Ensure the ROI contains fluorescent signal
    \item Try ``No background correction'' option
    \item Verify you're on the correct image/channel
\end{itemize}

\subsection{``No ROI selected''}

\textbf{Cause:} You clicked OK without drawing a selection.

\textbf{Solution:}
\begin{itemize}
    \item Draw the ROI first, then click OK
    \item For ``Use existing ROI'': ensure an ROI exists before running the script
\end{itemize}

\subsection{``Coordinates out of image bounds''}

\textbf{Cause:} Typed coordinates exceed image dimensions.

\textbf{Solution:}
\begin{itemize}
    \item Check image size: \texttt{Image > Properties}
    \item Coordinates are 0-indexed (0 to width-1, 0 to height-1)
\end{itemize}

\subsection{Results show NaN or Infinity}

\textbf{Cause:} Division by zero ($F_0 = 0$) or invalid measurements.

\textbf{Solution:}
\begin{itemize}
    \item Background point might be inside a bright region
    \item Try ``Background subtraction only'' method
    \item Check for corrupted frames in the stack
\end{itemize}

\subsection{Unexpected normalized values}

\textbf{Cause:} Intensity changes not matching expectations.

\textbf{Solution:}
\begin{itemize}
    \item Verify background point is in a non-fluorescent region
    \item Check for photobleaching during setup
    \item Ensure first frame is representative baseline
\end{itemize}

%==============================================================================
\section{Example Workflow}
%==============================================================================

\begin{enumerate}
    \item Open stack: \texttt{File > Open > my\_timelapse.tif}
    \item Run script: \texttt{Plugins > Macros > Run > Intensity\_measure.js}
    \item Normalization: Select ``F/F0 (Normalize to first frame)''
    \item Background: Click to select $\rightarrow$ Click on dark region $\rightarrow$ OK
    \item ROI: Freehand $\rightarrow$ Trace around cell $\rightarrow$ OK
    \item Results table appears
    \item Export: \texttt{File > Save As > intensity\_data.csv}
    \item Analyze/plot in your preferred software
\end{enumerate}

%==============================================================================
\section{File Structure}
%==============================================================================

\begin{lstlisting}[language={},numbers=none]
Ares/
    Align_Stack.js
    Crop_Stack_Subsets.js
    Intensity_measure.js
    Intensity_measure_documentation.tex
    README.md
    CHANGELOG.md
    CITATION.cff
    docs/
        intensity-measure.md
        align-stack.md
        stack-crop-subsets.md
\end{lstlisting}

%==============================================================================
\section{Changelog}
%==============================================================================

\begin{description}
    \item[v1.0] Initial release with basic rectangle ROI and point background
    \item[v1.1] Added coordinate input options
    \item[v1.2] Added F/F$_0$ photobleaching normalization
    \item[v1.3] Added multiple ROI selection tools (oval, freehand, polygon, brush)
    \item[v2.0] Added multiple normalization methods (F/F$_0$, $\Delta$F/F$_0$, background subtraction, percent max, Z-score, raw)
    \item[v2.1] Corrected CTCF area handling for non-rectangular ROIs and updated repository documentation structure
\end{description}

\vfill

\begin{center}
\rule{0.5\textwidth}{0.4pt}\\[0.5cm]
\textbf{Author:} Radmir Sarsenov\\[0.3cm]
\textbf{License:} MIT License\\[0.3cm]
\small Copyright \textcopyright{} 2026 Radmir Sarsenov
\end{center}

\end{document}
